\clearpage\pagestyle{otMain}\setlength{\otBlocEspacement}{2.5mm}\fontsize{11}{15}\selectfont\corrigesChapitre{8}{Grandeurs composées et unités}
\begin{blocLivre}{cor-8-1}\begin{corrige}{1}{Une durée}\demarche{Utiliser une heure de 60 minutes.}
$18$ min $=18/60=0{,}3$ h, donc la durée vaut $\res{2{,}3\ \mathrm h}$. En minutes : $2\times60+18=\res{138}$ min.
\end{corrige}\end{blocLivre}
\begin{blocLivre}{cor-8-2}\begin{corrige}{2}{Une vitesse moyenne}\demarche{Exprimer la durée dans l’unité du quotient.}
$2$ h $30$ min $=2{,}5$ h. La vitesse moyenne est $45/2{,}5=\res{18\ \mathrm{km/h}}$.
\end{corrige}\end{blocLivre}
\begin{blocLivre}{cor-8-3}\begin{corrige}{3}{Une conversion de vitesse}\demarche{Convertir distance et durée.}
$72$ km correspondent à $72\,000$ m en 3600 s. Donc $72\,000/3600=\res{20\ \mathrm{m/s}}$. On retrouve $72/3{,}6=20$.
\end{corrige}\end{blocLivre}
\begin{blocLivre}{cor-8-4}\begin{corrige}{4}{Un débit}\demarche{Diviser le volume par le temps.}
Le débit vaut volume sur durée : $q=18/45=\res{0{,}4\ \mathrm{L/s}}$. Cela signifie 0,4 litre en une seconde dans le modèle à débit constant.
\end{corrige}\end{blocLivre}
\begin{blocLivre}{cor-8-5}\begin{corrige}{5}{Un volume et une contenance}\demarche{Utiliser les correspondances entre volume et contenance.}
$0{,}35\times1000=\res{350}$ L. Comme 1 mL = 1 cm$^3$, $250$ mL $=\res{250\ \mathrm{cm}^3}$.
\end{corrige}\end{blocLivre}
\begin{blocLivre}{cor-8-6}\begin{corrige}{6}{Distance et temps}\demarche{Multiplier la vitesse par la durée dans la bonne unité.}
$35$ min $=35/60$ h. $d=v\times t=12\times35/60=\res7$ km. Une durée inférieure à une heure donne ici une distance inférieure à 12 km.
\end{corrige}\end{blocLivre}
\begin{blocLivre}{cor-8-7}\begin{corrige}{7}{Remplir une cuve}\demarche{Diviser le volume par le débit.}
$t=V/q=450/15=\res{30}$ min. Contrôle : $15\times30=450$ L.
\end{corrige}\end{blocLivre}
\begin{blocLivre}{cor-8-8}\begin{corrige}{8}{Changer les deux unités}\demarche{Convertir séparément le numérateur et le temps de référence.}
On convertit d’abord : $0{,}012$ m$^3$/s $=12$ L/s. En une minute, il s’écoule $12\times60=\res{720\ \mathrm{L/min}}$. Le temps est au dénominateur : une minute contient davantage de volume qu’une seconde.
\end{corrige}\end{blocLivre}
\begin{blocLivre}{cor-8-9}\begin{corrige}{9}{Une concentration}\demarche{Convertir le volume avant le quotient.}
Le volume final vaut $0{,}300$ L. $c=m/V=18/0{,}3=\res{60\ \mathrm{g/L}}$. Le volume utilisé est celui de la solution finale, tel qu’il est donné.
\end{corrige}\end{blocLivre}
\begin{blocLivre}{cor-8-10}\begin{corrige}{10}{Une énergie}\demarche{Associer chaque unité à la formule.}
$45$ min $=0{,}75$ h. Donc $E=1{,}2\times0{,}75=\res{0{,}9\ \mathrm{kWh}}$. La durée doit être en heures pour obtenir des kWh.
\end{corrige}\end{blocLivre}
\begin{blocLivre}{cor-8-11}\begin{corrige}{11}{Aller et retour}\demarche{Diviser la distance totale par la durée totale.}
L’aller dure $30/15=2$ h et le retour $30/30=1$ h. Pour 60 km en 3 h, la vitesse moyenne vaut $\res{20\ \mathrm{km/h}}$. La moyenne simple 22,5 ne tient pas compte des durées différentes.
\end{corrige}\end{blocLivre}
\begin{blocLivre}{cor-8-12}\begin{corrige}{12}{Une masse volumique}\demarche{Calculer le quotient puis réutiliser le coefficient.}
$\rho=108/40=\res{2{,}7\ \mathrm{g/cm}^3}$. À masse volumique identique, $m=\rho V=2{,}7\times65=\res{175{,}5}$ g.
\end{corrige}\end{blocLivre}
\begin{blocLivre}{cor-8-13}\begin{corrige}{13}{Deux arrivées d’eau}\demarche{Additionner les volumes fournis en une minute.}
Le volume vaut 1200 L. Les débits s’ajoutent : $8+12=20$ L/min. Donc $t=1200/20=\res{60}$ min, soit 1 h. Additionner les durées séparées serait incorrect.
\end{corrige}\end{blocLivre}
\begin{blocLivre}{cor-8-14}\begin{corrige}{14}{Une pause dans le trajet}\demarche{Inclure tous les temps du bilan demandé.}
La marche dure $9/4{,}5=2$ h. La durée totale est 2,5 h. La vitesse moyenne vaut $9/2{,}5=\res{3{,}6\ \mathrm{km/h}}$. Elle est inférieure à la vitesse de marche car la pause ajoute du temps sans ajouter de distance.
\end{corrige}\end{blocLivre}
\begin{blocLivre}{cor-8-15}\begin{corrige}{15}{Une comparaison de débits}\demarche{Comparer dans une même unité puis calculer les durées.}
B débite $0{,}018\times1000\times60=\res{1080}$ L/min, donc plus que A. A met $5400/900=6$ min ; B met $5400/1080=5$ min. Le gain est $\res1$ min.
\end{corrige}\end{blocLivre}
